% ===================== TÓM TẮT ========================================
\chapter*{Tóm tắt}
\addcontentsline{toc}{chapter}{Tóm tắt}

Học máy tự động (\textit{Automated Machine Learning}, AutoML) giúp tự động hoá
quá trình lựa chọn thuật toán và tinh chỉnh siêu tham số, nhờ đó rút ngắn đáng
kể công sức xây dựng mô hình. Tuy nhiên, việc so sánh các thư viện AutoML trong
thực tế thường thiếu công bằng và khó tái lập: mỗi thư viện được chạy với ngân
sách thời gian, cách chia dữ liệu, độ đo và tài nguyên khác nhau, đồng thời các
lần chạy thất bại thường bị bỏ qua thay vì được ghi nhận. Điều này khiến kết quả
so sánh dễ gây hiểu nhầm và khó kiểm chứng.

Trong đồ án này, nhóm tái hiện và mở rộng phương pháp đánh giá công bằng theo
chuẩn \textit{AutoML Benchmark} (AMLB) của Gijsbers và cộng sự, áp dụng cùng một
giao thức thống nhất (cùng cách chia fold, cùng ngân sách thời gian, cùng độ đo
và cùng tài nguyên) cho mọi thư viện trên mỗi tập dữ liệu, đồng thời ghi nhận
đầy đủ các lần chạy thất bại. Nhóm xây dựng một hệ thống \textit{AutoML Bench
Console} cho phép nạp dữ liệu (tải lên tệp CSV, lấy từ OpenML hoặc nhập liên kết
Kaggle công khai), tích hợp thư viện AutoML dưới dạng ảnh Docker, chạy đánh giá
và trực quan hoá kết quả, với dữ liệu được lưu trữ trên PostgreSQL và MinIO.

Nhóm tiến hành thực nghiệm so sánh ba thư viện AutoML tiêu biểu là FLAML,
H2O~AutoML và AutoGluon trên các tập dữ liệu phân lớp và hồi quy, sử dụng các độ
đo AUC, log-loss và RMSE, và xếp hạng các thư viện theo thứ hạng trung bình. Kết
quả sơ bộ cho thấy [tóm tắt phát hiện chính sau khi hoàn tất lần chạy đầy đủ], và
làm rõ sự đánh đổi giữa chất lượng mô hình và chi phí thời gian/tài nguyên giữa
các thư viện.

\vspace{0.5cm}
\noindent\textbf{Từ khoá:} học máy tự động (AutoML), đánh giá công bằng
(benchmark), khả năng tái lập, FLAML, H2O AutoML, AutoGluon.
